\begin{register}{H}{GPIO\_BT\_SELECT\_REG}{0x{}0000}\label{regdesc:GPIOBTSELECTREG}
\regfield{GPIO\_BT\_SEL}{32}{0}{{0x{}000000}}%
\reglabel{Reset}\regnewline%
\begin{regdesc}\begin{reglist}
\label{fielddesc:GPIOBTSEL}\item [GPIO\_BT\_SEL] 保留。(R/W)
\end{reglist}\end{regdesc}
\end{register}


\begin{register}{H}{GPIO\_OUT\_REG}{0x{}0004}\label{regdesc:GPIOOUTREG}
\regfield{GPIO\_OUT\_DATA\_ORIG}{32}{0}{{0x{}000000}}%
\reglabel{Reset}\regnewline%
\begin{regdesc}\begin{reglist}
\label{fielddesc:GPIOOUTDATAORIG}\item [GPIO\_OUT\_DATA\_ORIG] 简单 GPIO 输出模式下，GPIO0 \verb+~+ 21 和 GPIO26 \verb+~+31 的输出值。bit0 \verb+~+ bit21 的值分别对应 GPIO0 \verb+~+ GPIO21 的输出值；bit26 \verb+~+ bit31 的值分别对应 GPIO26 \verb+~+ GPIO31 的输出值。bit22 \verb+~+ bit25 无效。(R/W)
\end{reglist}\end{regdesc}
\end{register}


\begin{register}{H}{GPIO\_OUT\_W1TS\_REG}{0x{}0008}\label{regdesc:GPIOOUTW1TSREG}
\regfield{GPIO\_OUT\_W1TS}{32}{0}{{0x{}000000}}%
\reglabel{Reset}\regnewline%
\begin{regdesc}\begin{reglist}
\label{fielddesc:GPIOOUTW1TS}\item [GPIO\_OUT\_W1TS] GPIO0 \verb+~+ 31 输出置位寄存器。每一位置 1， \hyperref[regdesc:GPIOOUTREG]{GPIO\_OUT\_REG} 中的相应位也置 1。注：推荐使用此寄存器来置位 \hyperref[regdesc:GPIOOUTREG]{GPIO\_OUT\_REG}。(WO)
\end{reglist}\end{regdesc}
\end{register}


\begin{register}{H}{GPIO\_OUT\_W1TC\_REG}{0x{}000C}\label{regdesc:GPIOOUTW1TCREG}
\regfield{GPIO\_OUT\_W1TC}{32}{0}{{0x{}000000}}%
\reglabel{Reset}\regnewline%
\begin{regdesc}\begin{reglist}
\label{fielddesc:GPIOOUTW1TC}\item [GPIO\_OUT\_W1TC] GPIO0 \verb+~+ 31 输出清零寄存器。每一位置 1，则 \hyperref[regdesc:GPIOOUTREG]{GPIO\_OUT\_REG} 中的相应位会清零。注：推荐使用此寄存器来清零 \hyperref[regdesc:GPIOOUTREG]{GPIO\_OUT\_REG}。(WO)
\end{reglist}\end{regdesc}
\end{register}


\begin{register}{H}{GPIO\_OUT1\_REG}{0x{}0010}\label{regdesc:GPIOOUT1REG}
\regfield{(reserved)}{10}{22}{0000000000}%
\regfield{GPIO\_OUT1\_DATA\_ORIG}{22}{0}{{0x{}0000}}%
\reglabel{Reset}\regnewline%
\begin{regdesc}\begin{reglist}
\label{fielddesc:GPIOOUT1DATAORIG}\item [GPIO\_OUT1\_DATA\_ORIG] 简单 GPIO 输出模式下，GPIO32 \verb+~+ 48 的输出值。bit0 \verb+~+ bit16 的值分别对应 GPIO32 \verb+~+ GPIO48 的输出值。bit17 \verb+~+ bit21 无效。(R/W)
\end{reglist}\end{regdesc}
\end{register}


\begin{register}{H}{GPIO\_OUT1\_W1TS\_REG}{0x{}0014}\label{regdesc:GPIOOUT1W1TSREG}
\regfield{(reserved)}{10}{22}{0000000000}%
\regfield{GPIO\_OUT1\_W1TS}{22}{0}{{0x{}0000}}%
\reglabel{Reset}\regnewline%
\begin{regdesc}\begin{reglist}
\label{fielddesc:GPIOOUT1W1TS}\item [GPIO\_OUT1\_W1TS] GPIO32 \verb+~+ 48 输出置位寄存器。每一位置 1，则 \hyperref[regdesc:GPIOOUT1REG]{GPIO\_OUT1\_REG} 中的相应位也置 1。注：推荐使用此寄存器来置位 \hyperref[regdesc:GPIOOUT1REG]{GPIO\_OUT1\_REG}。(WO)
\end{reglist}\end{regdesc}
\end{register}


\begin{register}{H}{GPIO\_OUT1\_W1TC\_REG}{0x{}0018}\label{regdesc:GPIOOUT1W1TCREG}
\regfield{(reserved)}{10}{22}{0000000000}%
\regfield{GPIO\_OUT1\_W1TC}{22}{0}{{0x{}0000}}%
\reglabel{Reset}\regnewline%
\begin{regdesc}\begin{reglist}
\label{fielddesc:GPIOOUT1W1TC}\item [GPIO\_OUT1\_W1TC] GPIO32 \verb+~+ 48 输出清零寄存器。每一位置 1，则 \hyperref[regdesc:GPIOOUT1REG]{GPIO\_OUT1\_REG} 中的相应位会清零。注：推荐使用此寄存器来清零 \hyperref[regdesc:GPIOOUT1REG]{GPIO\_OUT1\_REG}。(WO)
\end{reglist}\end{regdesc}
\end{register}


\begin{register}{H}{GPIO\_SDIO\_SELECT\_REG}{0x{}001C}\label{regdesc:GPIOSDIOSELECTREG}
\regfield{(reserved)}{24}{8}{000000000000000000000000}%
\regfield{GPIO\_SDIO\_SEL}{8}{0}{{0x{}0}}%
\reglabel{Reset}\regnewline%
\begin{regdesc}\begin{reglist}
\label{fielddesc:GPIOSDIOSEL}\item [GPIO\_SDIO\_SEL] 保留。(R/W)
\end{reglist}\end{regdesc}
\end{register}


\begin{register}{H}{GPIO\_ENABLE\_REG}{0x{}0020}\label{regdesc:GPIOENABLEREG}
\regfield{GPIO\_ENABLE\_DATA}{32}{0}{{0x{}000000}}%
\reglabel{Reset}\regnewline%
\begin{regdesc}\begin{reglist}
\label{fielddesc:GPIOENABLEDATA}\item [GPIO\_ENABLE\_DATA] GPIO0 \verb+~+ 31 输出使能寄存器。(R/W)
\end{reglist}\end{regdesc}
\end{register}


\begin{register}{H}{GPIO\_ENABLE\_W1TS\_REG}{0x{}0024}\label{regdesc:GPIOENABLEW1TSREG}
\regfield{GPIO\_ENABLE\_W1TS}{32}{0}{{0x{}000000}}%
\reglabel{Reset}\regnewline%
\begin{regdesc}\begin{reglist}
\label{fielddesc:GPIOENABLEW1TS}\item [GPIO\_ENABLE\_W1TS] GPIO0 \verb+~+ 31 输出使能置位寄存器。每一位置 1，则 \hyperref[regdesc:GPIOENABLEREG]{GPIO\_ENABLE\_REG} 中的相应位也置 1。注：推荐使用此寄存器来置位 \hyperref[regdesc:GPIOENABLEREG]{GPIO\_ENABLE\_REG}。(WO)
\end{reglist}\end{regdesc}
\end{register}


\begin{register}{H}{GPIO\_ENABLE\_W1TC\_REG}{0x{}0028}\label{regdesc:GPIOENABLEW1TCREG}
\regfield{GPIO\_ENABLE\_W1TC}{32}{0}{{0x{}000000}}%
\reglabel{Reset}\regnewline%
\begin{regdesc}\begin{reglist}
\label{fielddesc:GPIOENABLEW1TC}\item [GPIO\_ENABLE\_W1TC] GPIO0 \verb+~+ 31 输出使能清零寄存器。每一位置 1，则 \hyperref[regdesc:GPIOENABLEREG]{GPIO\_ENABLE\_REG} 中的相应位会清零。注：推荐使用此寄存器清零 \hyperref[regdesc:GPIOENABLEREG]{GPIO\_ENABLE\_REG}。(WO)
\end{reglist}\end{regdesc}
\end{register}


\begin{register}{H}{GPIO\_ENABLE1\_REG}{0x{}002C}\label{regdesc:GPIOENABLE1REG}
\regfield{(reserved)}{10}{22}{0000000000}%
\regfield{GPIO\_ENABLE1\_DATA}{22}{0}{{0x{}0000}}%
\reglabel{Reset}\regnewline%
\begin{regdesc}\begin{reglist}
\label{fielddesc:GPIOENABLE1DATA}\item [GPIO\_ENABLE1\_DATA] GPIO32 \verb+~+ 48 输出使能寄存器。(R/W)
\end{reglist}\end{regdesc}
\end{register}


\begin{register}{H}{GPIO\_ENABLE1\_W1TS\_REG}{0x{}0030}\label{regdesc:GPIOENABLE1W1TSREG}
\regfield{(reserved)}{10}{22}{0000000000}%
\regfield{GPIO\_ENABLE1\_W1TS}{22}{0}{{0x{}0000}}%
\reglabel{Reset}\regnewline%
\begin{regdesc}\begin{reglist}
\label{fielddesc:GPIOENABLE1W1TS}\item [GPIO\_ENABLE1\_W1TS] GPIO32 \verb+~+ 48 输出使能置位寄存器。每一位置 1，则 \hyperref[regdesc:GPIOENABLE1REG]{GPIO\_ENABLE1\_REG} 中的相应位也置 1。注：推荐使用此寄存器来置位 \hyperref[regdesc:GPIOENABLE1REG]{GPIO\_ENABLE1\_REG}。(WO)
\end{reglist}\end{regdesc}
\end{register}


\begin{register}{H}{GPIO\_ENABLE1\_W1TC\_REG}{0x{}0034}\label{regdesc:GPIOENABLE1W1TCREG}
\regfield{(reserved)}{10}{22}{0000000000}%
\regfield{GPIO\_ENABLE1\_W1TC}{22}{0}{{0x{}0000}}%
\reglabel{Reset}\regnewline%
\begin{regdesc}\begin{reglist}
\label{fielddesc:GPIOENABLE1W1TC}\item [GPIO\_ENABLE1\_W1TC] GPIO32 \verb+~+ 48 输出使能清零寄存器。每一位置 1，则 \hyperref[regdesc:GPIOENABLE1REG]{GPIO\_ENABLE1\_REG} 中的相应位会清零。注：推荐使用此寄存器清零 \hyperref[regdesc:GPIOENABLE1REG]{GPIO\_ENABLE1\_REG}。(WO)
\end{reglist}\end{regdesc}
\end{register}


\begin{register}{H}{GPIO\_STRAP\_REG}{0x{}0038}\label{regdesc:GPIOSTRAPREG}
\regfield{(reserved)}{26}{6}{00000000000000000000000000}%
\regfield{GPIO\_STRAPPING}{6}{0}{{0}}%
\reglabel{Reset}\regnewline%
\begin{regdesc}\begin{reglist}
\label{fielddesc:GPIOSTRAPPING}\item [GPIO\_STRAPPING] 表示对应 GPIO 管脚的 Strapping 值。\\
bit0 \verb+~+ bit1：保留\\
bit2：GPIO46\\
bit3：GPIO0\\
bit4：GPIO45\\
bit5：GPIO3\\(RO)
\end{reglist}\end{regdesc}
\end{register}


\begin{register}{H}{GPIO\_IN\_REG}{0x{}003C}\label{regdesc:GPIOINREG}
\regfield{GPIO\_IN\_DATA\_NEXT}{32}{0}{{0}}%
\reglabel{Reset}\regnewline%
\begin{regdesc}\begin{reglist}
\label{fielddesc:GPIOINDATANEXT}\item [GPIO\_IN\_DATA\_NEXT] GPIO0 \verb+~+ 31 输入值。每一位代表一个管脚的片外输入值。比如片外引脚为高电平，则此位应为 1；片外引脚为低电平，此位应为 0。(RO)
\end{reglist}\end{regdesc}
\end{register}


\begin{register}{H}{GPIO\_IN1\_REG}{0x{}0040}\label{regdesc:GPIOIN1REG}
\regfield{(reserved)}{10}{22}{0000000000}%
\regfield{GPIO\_IN\_DATA1\_NEXT}{22}{0}{{0}}%
\reglabel{Reset}\regnewline%
\begin{regdesc}\begin{reglist}
\label{fielddesc:GPIOINDATA1NEXT}\item [GPIO\_IN\_DATA1\_NEXT] GPIO32 \verb+~+ 48 输入值。每一位代表一个管脚的片外输入值。(RO)
\end{reglist}\end{regdesc}
\end{register}


\begin{register}{H}{GPIO\_PIN\regindex{n}\_REG (\regindex{n}: 0-48)}{0x{}0074+0x{}4*\regindex{n}}\label{regdesc:GPIOPINNREG}
\regfield{(reserved)}{14}{18}{00000000000000}%
\regfield{GPIO\_PIN\regindex{n}\_INT\_ENA}{5}{13}{{0x{}0}}%
\regfield{GPIO\_PIN\regindex{n}\_CONFIG}{2}{11}{{0x{}0}}%
\regfield{GPIO\_PIN\regindex{n}\_WAKEUP\_ENABLE}{1}{10}{{0}}%
\regfield{GPIO\_PIN\regindex{n}\_INT\_TYPE}{3}{7}{{0x{}0}}%
\regfield{(reserved)}{2}{5}{00}%
\regfield{GPIO\_PIN\regindex{n}\_SYNC1\_BYPASS}{2}{3}{{0x{}0}}%
\regfield{GPIO\_PIN\regindex{n}\_PAD\_DRIVER}{1}{2}{{0}}%
\regfield{GPIO\_PIN\regindex{n}\_SYNC2\_BYPASS}{2}{0}{{0x{}0}}%
\reglabel{Reset}\regnewline%
\begin{regdesc}\begin{reglist}
\label{fielddesc:GPIOPINNSYNC2BYPASS}\item [GPIO\_PIN\regindex{n}\_SYNC2\_BYPASS] 使能 GPIO 输入信号第二拍为 APB 时钟上升沿或下降沿同步。0：关闭同步；1：下降沿同步；2 或 3：上升沿同步。(R/W)
\label{fielddesc:GPIOPINNPADDRIVER}\item [GPIO\_PIN\regindex{n}\_PAD\_DRIVER] 管脚驱动选择。0：正常输出；1：开漏输出。(R/W)
\label{fielddesc:GPIOPINNSYNC1BYPASS}\item [GPIO\_PIN\regindex{n}\_SYNC1\_BYPASS] 使能 GPIO 输入信号第一拍为 APB 时钟上升沿或下降沿同步。0：关闭同步；1：下降沿同步；2 或 3：上升沿同步。(R/W)
\label{fielddesc:GPIOPINNINTTYPE}\item [GPIO\_PIN\regindex{n}\_INT\_TYPE] 中断类型选择。(R/W)

    0：禁用 GPIO 中断 \\
    1：上升沿触发 \\
    2：下降沿触发 \\
    3：任一沿触发 \\
    4：低电平触发 \\
    5：高电平触发 \\

\label{fielddesc:GPIOPINNWAKEUPENABLE}\item [GPIO\_PIN\regindex{n}\_WAKEUP\_ENABLE] 使能 GPIO 唤醒，仅能将 CPU 从 Light-sleep 模式唤醒。(R/W)
\label{fielddesc:GPIOPINNCONFIG}\item [GPIO\_PIN\regindex{n}\_CONFIG] 保留。(R/W)
\label{fielddesc:GPIOPINNINTENA}\item [GPIO\_PIN\regindex{n}\_INT\_ENA] 中断使能位。bit13：使能 CPU 中断；bit14：使能 CPU 非屏蔽中断。(R/W)
\end{reglist}\end{regdesc}
\end{register}


\begin{register}{H}{GPIO\_FUNC\regindex{y}\_IN\_SEL\_CFG\_REG (\regindex{y}: 0-255)}{0x{}0154+0x{}4*\regindex{y}}\label{regdesc:GPIOFUNCNINSELCFGREG}
\regfield{(reserved)}{24}{8}{000000000000000000000000}%
\regfield{GPIO\_SIG\regindex{y}\_IN\_SEL}{1}{7}{{0}}%
\regfield{GPIO\_FUNC\regindex{y}\_IN\_INV\_SEL}{1}{6}{{0}}%
\regfield{GPIO\_FUNC\regindex{y}\_IN\_SEL}{6}{0}{{0x{}0}}%
\reglabel{Reset}\regnewline%
\begin{regdesc}\begin{reglist}
\label{fielddesc:GPIOFUNCNINSEL}\item [GPIO\_FUNC\regindex{y}\_IN\_SEL] 外设输入信号 \regindex{Y} 的选择控制位。此位选择 1 个 GPIO 交换矩阵输入管脚与信号连接；或者选择 0x{}38，则输入信号恒为高电平；或者选择 0x{}3C，则输入信号恒为低电平。(R/W)
\label{fielddesc:GPIOFUNCNININVSEL}\item [GPIO\_FUNC\regindex{y}\_IN\_INV\_SEL] 反转输入值。1：反转；0：不反转。(R/W)
\label{fielddesc:GPIOSIGNINSEL}\item [GPIO\_SIG\regindex{y}\_IN\_SEL] 旁路 GPIO 交换矩阵。1：通过 GPIO 交换矩阵；0：直接通过 IO MUX 连接信号与外设。(R/W)
\end{reglist}\end{regdesc}
\end{register}


\begin{register}{H}{GPIO\_FUNC\regindex{x}\_OUT\_SEL\_CFG\_REG (\regindex{x}: 0-48)}{0x{}0554+0x{}4*\regindex{x}}\label{regdesc:GPIOFUNCNOUTSELCFGREG}
\regfield{(reserved)}{20}{12}{00000000000000000000}%
\regfield{GPIO\_FUNC\regindex{x}\_OEN\_INV\_SEL}{1}{11}{{0}}%
\regfield{GPIO\_FUNC\regindex{x}\_OEN\_SEL}{1}{10}{{0}}%
\regfield{GPIO\_FUNC\regindex{x}\_OUT\_INV\_SEL}{1}{9}{{0}}%
\regfield{GPIO\_FUNC\regindex{x}\_OUT\_SEL}{9}{0}{{0x{}100}}%
\reglabel{Reset}\regnewline%
\begin{regdesc}\begin{reglist}
\label{fielddesc:GPIOFUNCNOUTSEL}\item [GPIO\_FUNC\regindex{x}\_OUT\_SEL] GPIO 管脚 \regindex{X} 的输出信号选择控制位。值为 \regindex{y} (0<=\regindex{y}<256) 连接外设输出 \regindex{y} 与 GPIO 输出 \regindex{x}。值为 256 选择 \hyperref[regdesc:GPIOOUTREG]{GPIO\_OUT\_REG}/\hyperref[regdesc:GPIOOUT1REG]{GPIO\_OUT1\_REG}[\regindex{x}] 和 \hyperref[regdesc:GPIOENABLEREG]{GPIO\_ENABLE\_REG}/\hyperref[regdesc:GPIOENABLE1REG]{GPIO\_ENABLE1\_REG} [\regindex{x}] 作为输出值和输出使能。(R/W)
\label{fielddesc:GPIOFUNCNOUTINVSEL}\item [GPIO\_FUNC\regindex{x}\_OUT\_INV\_SEL] 0：不反转输出值；1：反转输出值。(R/W)
\label{fielddesc:GPIOFUNCNOENSEL}\item [GPIO\_FUNC\regindex{x}\_OEN\_SEL] 0：采用外设的输出使能信号；1：强制使用 \hyperref[regdesc:GPIOENABLEREG]{GPIO\_ENABLE\_REG}[\regindex{x}] 用作输出使能信号。(R/W)
\label{fielddesc:GPIOFUNCNOENINVSEL}\item [GPIO\_FUNC\regindex{x}\_OEN\_INV\_SEL] 0：不反转输出使能信号；1：反转输出使能信号。(R/W)
\end{reglist}\end{regdesc}
\end{register}


\begin{register}{H}{GPIO\_CLOCK\_GATE\_REG}{0x{}062C}\label{regdesc:GPIOCLOCKGATEREG}
\regfield{(reserved)}{31}{1}{0000000000000000000000000000000}%
\regfield{GPIO\_CLK\_EN}{1}{0}{{1}}%
\reglabel{Reset}\regnewline%
\begin{regdesc}\begin{reglist}
\label{fielddesc:GPIOCLKEN}\item [GPIO\_CLK\_EN] 时钟门控使能。此位置 1，则时钟自由运转。(R/W)
\end{reglist}\end{regdesc}
\end{register}


\begin{register}{H}{GPIO\_STATUS\_REG}{0x{}0044}\label{regdesc:GPIOSTATUSREG}
\regfield{GPIO\_STATUS\_INTERRUPT}{32}{0}{{0x{}000000}}%
\reglabel{Reset}\regnewline%
\begin{regdesc}\begin{reglist}
\label{fielddesc:GPIOSTATUSINTERRUPT}\item [GPIO\_STATUS\_INTERRUPT] GPIO0 \verb+~+ 31 中断状态寄存器。(R/W)
\end{reglist}\end{regdesc}
\end{register}


\begin{register}{H}{GPIO\_STATUS1\_REG}{0x{}0050}\label{regdesc:GPIOSTATUS1REG}
\regfield{(reserved)}{10}{22}{0000000000}%
\regfield{GPIO\_STATUS1\_INTERRUPT}{22}{0}{{0x{}0000}}%
\reglabel{Reset}\regnewline%
\begin{regdesc}\begin{reglist}
\label{fielddesc:GPIOSTATUS1INTERRUPT}\item [GPIO\_STATUS1\_INTERRUPT] GPIO32 \verb+~+ 48 中断状态寄存器。(R/W)
\end{reglist}\end{regdesc}
\end{register}


\begin{register}{H}{GPIO\_CPU\_INT\_REG}{0x{}005C}\label{regdesc:GPIOPCPUINTREG}
\regfield{GPIO\_CPU\_INT}{32}{0}{{0x{}000000}}%
\reglabel{Reset}\regnewline%
\begin{regdesc}\begin{reglist}
\label{fielddesc:GPIOPROCPUINT}\item [GPIO\_CPU\_INT] GPIO0 \verb+~+ 31 CPU 中断状态。如果 \hyperref[regdesc:GPIOPINNREG]{GPIO\_PIN\regindex{n}\_REG} 中 bit13 高电平有效，即使能 CPU 中断，则此寄存器所示的中断状态应与 \hyperref[regdesc:GPIOSTATUSREG]{GPIO\_STATUS\_REG} 中相应 bit 的中断状态一致。(RO)
\end{reglist}\end{regdesc}
\end{register}


\begin{register}{H}{GPIO\_CPU\_NMI\_INT\_REG}{0x{}0060}\label{regdesc:GPIOPCPUNMIINTREG}
\regfield{GPIO\_CPU\_NMI\_INT}{32}{0}{{0x{}000000}}%
\reglabel{Reset}\regnewline%
\begin{regdesc}\begin{reglist}
\label{fielddesc:GPIOPROCPUNMIINT}\item [GPIO\_CPU\_NMI\_INT] GPIO0 \verb+~+ 31 CPU 非屏蔽中断状态寄存器。如果 \hyperref[regdesc:GPIOPINNREG]{GPIO\_PIN\regindex{n}\_REG} 中 bit14 高电平有效，即使能 CPU 非屏蔽中断，则此寄存器所示的中断状态应与 \hyperref[regdesc:GPIOSTATUSREG]{GPIO\_STATUS\_REG} 中相应 bit 的中断状态一致。(RO)
\end{reglist}\end{regdesc}
\end{register}


\begin{register}{H}{GPIO\_CPU\_INT1\_REG}{0x{}0068}\label{regdesc:GPIOPCPUINT1REG}
\regfield{(reserved)}{10}{22}{0000000000}%
\regfield{GPIO\_CPU1\_INT}{22}{0}{{0x{}0000}}%
\reglabel{Reset}\regnewline%
\begin{regdesc}\begin{reglist}
\label{fielddesc:GPIOPROCPU1INT}\item [GPIO\_CPU1\_INT] GPIO32 \verb+~+ 48 CPU 中断状态寄存器。如果 \hyperref[regdesc:GPIOPINNREG]{GPIO\_PIN\regindex{n}\_REG} 中 bit13 高电平有效，即使能 CPU 中断，则此寄存器所示的中断状态应与 \hyperref[regdesc:GPIOSTATUS1REG]{GPIO\_STATUS1\_REG} 中相应 bit 的中断状态一致。(RO)
\end{reglist}\end{regdesc}
\end{register}


\begin{register}{H}{GPIO\_CPU\_NMI\_INT1\_REG}{0x{}006C}\label{regdesc:GPIOPCPUNMIINT1REG}
\regfield{(reserved)}{10}{22}{0000000000}%
\regfield{GPIO\_CPU\_NMI1\_INT}{22}{0}{{0x{}0000}}%
\reglabel{Reset}\regnewline%
\begin{regdesc}\begin{reglist}
\label{fielddesc:GPIOPROCPUNMI1INT}\item [GPIO\_CPU\_NMI1\_INT] GPIO32 \verb+~+ 48 CPU 非屏蔽中断状态寄存器。如果 \hyperref[regdesc:GPIOPINNREG]{GPIO\_PIN\regindex{n}\_REG} 中 bit14 高电平有效，即使能 CPU 非屏蔽中断，则此寄存器所示的中断状态应与 \hyperref[regdesc:GPIOSTATUS1REG]{GPIO\_STATUS1\_REG} 中相应 bit 的中断状态一致。(RO)
\end{reglist}\end{regdesc}
\end{register}


\begin{register}{H}{GPIO\_STATUS\_W1TS\_REG}{0x{}0048}\label{regdesc:GPIOSTATUSW1TSREG}
\regfield{GPIO\_STATUS\_W1TS}{32}{0}{{0x{}000000}}%
\reglabel{Reset}\regnewline%
\begin{regdesc}\begin{reglist}
\label{fielddesc:GPIOSTATUSW1TS}\item [GPIO\_STATUS\_W1TS] GPIO0 \verb+~+ 31 中断状态置位寄存器。每位置 1，则  \hyperref[fielddesc:GPIOSTATUSINTERRUPT]{GPIO\_STATUS\_INTERRUPT} 中的相应位也置 1。注：推荐使用此寄存器来置位 \hyperref[fielddesc:GPIOSTATUSINTERRUPT]{GPIO\_STATUS\_INTERRUPT}。(WO)
\end{reglist}\end{regdesc}
\end{register}


\begin{register}{H}{GPIO\_STATUS\_W1TC\_REG}{0x{}004C}\label{regdesc:GPIOSTATUSW1TCREG}
\regfield{GPIO\_STATUS\_W1TC}{32}{0}{{0x{}000000}}%
\reglabel{Reset}\regnewline%
\begin{regdesc}\begin{reglist}
\label{fielddesc:GPIOSTATUSW1TC}\item [GPIO\_STATUS\_W1TC] GPIO0 \verb+~+ 31 中断状态清除寄存器。每一位置 1，则  \hyperref[fielddesc:GPIOSTATUSINTERRUPT]{GPIO\_STATUS\_INTERRUPT} 中的相应位也会清零。注：推荐使用此寄存器来清零 \hyperref[fielddesc:GPIOSTATUSINTERRUPT]{GPIO\_STATUS\_INTERRUPT}。(WO)
\end{reglist}\end{regdesc}
\end{register}


\begin{register}{H}{GPIO\_STATUS1\_W1TS\_REG}{0x{}0054}\label{regdesc:GPIOSTATUS1W1TSREG}
\regfield{(reserved)}{10}{22}{0000000000}%
\regfield{GPIO\_STATUS1\_W1TS}{22}{0}{{0x{}0000}}%
\reglabel{Reset}\regnewline%
\begin{regdesc}\begin{reglist}
\label{fielddesc:GPIOSTATUS1W1TS}\item [GPIO\_STATUS1\_W1TS] GPIO32 \verb+~+ 48 中断状态置位寄存器。每一位置 1， \hyperref[fielddesc:GPIOSTATUS1INTERRUPT]{GPIO\_STATUS\_INTERRUPT1} 中相应位也置 1。注：推荐使用此寄存器来置位 \hyperref[fielddesc:GPIOSTATUS1INTERRUPT]{GPIO\_STATUS\_INTERRUPT1}。(WO)
\end{reglist}\end{regdesc}
\end{register}


\begin{register}{H}{GPIO\_STATUS1\_W1TC\_REG}{0x{}0058}\label{regdesc:GPIOSTATUS1W1TCREG}
\regfield{(reserved)}{10}{22}{0000000000}%
\regfield{GPIO\_STATUS1\_W1TC}{22}{0}{{0x{}0000}}%
\reglabel{Reset}\regnewline%
\begin{regdesc}\begin{reglist}
\label{fielddesc:GPIOSTATUS1W1TC}\item [GPIO\_STATUS1\_W1TC] GPIO32 \verb+~+ 48 中断状态清除寄存器。每一位置 1，则  \hyperref[fielddesc:GPIOSTATUS1INTERRUPT]{GPIO\_STATUS\_INTERRUPT1} 中的相应位也将清零。注：推荐使用此寄存器来清零 \hyperref[fielddesc:GPIOSTATUS1INTERRUPT]{GPIO\_STATUS\_INTERRUPT1}。(WO)
\end{reglist}\end{regdesc}
\end{register}


\begin{register}{H}{GPIO\_STATUS\_NEXT\_REG}{0x{}014C}\label{regdesc:GPIOSTATUSNEXTREG}
\regfield{GPIO\_STATUS\_INTERRUPT\_NEXT}{32}{0}{{0x{}000000}}%
\reglabel{Reset}\regnewline%
\begin{regdesc}\begin{reglist}
\label{fielddesc:GPIOSTATUSINTERRUPTNEXT}\item [GPIO\_STATUS\_INTERRUPT\_NEXT] GPIO0 \verb+~+ 31 中断源信号，可以设置为上升沿中断、下降沿中断、电平敏感中断或任一沿中断。(RO)
\end{reglist}\end{regdesc}
\end{register}


\begin{register}{H}{GPIO\_STATUS\_NEXT1\_REG}{0x{}0150}\label{regdesc:GPIOSTATUSNEXT1REG}
\regfield{(reserved)}{10}{22}{0000000000}%
\regfield{GPIO\_STATUS1\_INTERRUPT\_NEXT}{22}{0}{{0x{}0000}}%
\reglabel{Reset}\regnewline%
\begin{regdesc}\begin{reglist}
\label{fielddesc:GPIOSTATUS1INTERRUPTNEXT}\item [GPIO\_STATUS1\_INTERRUPT\_NEXT] GPIO32 \verb+~+ 48 的中断源信号。(RO)
\end{reglist}\end{regdesc}
\end{register}


\begin{register}{H}{GPIO\_REG\_DATE\_REG}{0x{}06FC}\label{regdesc:GPIOREGDATEREG}
\regfield{(reserved)}{4}{28}{0000}%
\regfield{GPIO\_DATE}{28}{0}{{0x{}1907040}}%
\reglabel{Reset}\regnewline%
\begin{regdesc}\begin{reglist}
\label{fielddesc:GPIODATE}\item [GPIO\_DATE] 版本控制寄存器。(R/W)
\end{reglist}\end{regdesc}
\end{register}


